\documentclass[11pt]{article}
\usepackage{amsmath,amssymb,amsthm,hyperref}
\newtheorem{theorem}{Theorem}
\newcommand{\E}{\mathbb E}
\title{Converting a quantitative endpoint profile into a face bound}
\author{Research lemma: independently reviewed}
\date{30 September 2026}
\begin{document}\maketitle
This is a conditional quantitative strengthening of the fully reviewed
three-halves argument in \path{../power_per_die_lower_bound.tex}.
It does not assert the additional profile hypothesis without proof.

Let $p_j:[K]\to[0,1]$ be nondecreasing, with positive row weights $w$,
and suppose $\E_w\prod_{j\in S}p_j=1/(|S|+1)$ for every $S\subseteq[m]$.
Put $u_j=1-p_j$, $s=\sum_j u_j$, $\ell=\lfloor m/4\rfloor$,
$\alpha=\ell/m$, and
\[
 H=\ell(\max_j u_j-\min_j u_j).
\]
The finite endpoint-balance theorem gives $H\le C(s+1)$ and
\begin{equation}\label{laplace}
 \E_w(s+1)^k e^{-as}\le C_a^{k+1}k!/m
 \qquad(k\ge0)
\end{equation}
for every fixed $a>0$.

\begin{theorem}\label{main}
There are absolute constants $c,C>0$ with the following property.
Suppose, for some $1\le L\le m/256$ and $0\le\delta\le1$, that
\begin{equation}\label{extra}
 H\le\delta(s+1)\qquad\hbox{on all rows with }s\le L.
\end{equation}
Write $\rho=\delta+e^{-cL}$. Then
\[
 w_K\le C\max\{m^{-2},\,m^{-3/2}\sqrt\rho\}.
\]
For an equal-weight representation,
\[
 K\ge c\min\{m^2,\,m^{3/2}/\sqrt\rho\}.
\]
The analogous estimate for $w_1$ holds if the corresponding reflected
profile hypothesis is available.
\end{theorem}
\begin{proof}
Retain $\theta=1/16$, the success-subset mixture
$\mu=(\ell+1)w\overline A$, $A_B=\prod_{j\in B}p_j$, and the beta law
$\beta_\ell$ from the three-halves proof. For Laguerre sums
$Q_a(z)=\sum_{j=0}^d c_{aj}L_j(\theta z)$ write
$B_a=\sum|c_{aj}|$ and $B=B_1B_2$.
We first prove, for an absolute sufficiently small $c_0>0$ and
$1\le d\le c_0m$, the uniform estimate
\begin{equation}\label{comparison}
 |\E_\mu Q_1(\alpha s)Q_2(\alpha s)
       -\beta_\ell(Q_1Q_2)|
 \le CB\frac d m(\delta+e^{-cL}).
\end{equation}
Taking $c_0<3/8$ ensures that the degree is below the number of
remaining columns.

For a complement $R$ of size $N=m-\ell$, set $z_j=\ell u_j$ and
$v=N^{-1}\sum_{j\in R}z_j$. The exact polarization expansion,
the entire centered-coefficient estimate, and the variance-adapted
Laguerre derivative bound give a pointwise error at most
\[
 B e^{\theta v}\sum_{k=2}^{2d}
       \frac{(C\sqrt{dkH/m})^k}{k!}.
\]
Indeed $\sum_{j\in R}(z_j-v)^2\le NHv$, so all singular powers of
$v$ cancel. If $v=0$, the error is identically zero.
On $s\le L$, substitute $H\le\delta(s+1)$, multiply by the success
weight, and use \eqref{laplace}. The resulting sum is bounded by
\[
 CB\sum_{k=2}^{2d}(C\sqrt{d\delta/m})^k(k+1)
 \le CBd\delta/m.
\]
On $s>L$, reserve half of the exponential decay before applying
\eqref{laplace}: $e^{-as}\boldsymbol1_{s>L}
\le e^{-aL/2}e^{-as/2}$. The coarse bound $H\le C(s+1)$ then gives
\[
 CB e^{-cL}\sum_{k=2}^{2d}(C\sqrt{d/m})^k(k+1)
 \le CB(d/m)e^{-cL}.
\]
Both geometric estimates hold by choosing $c_0$ absolutely small.

For completeness the success-subset averaging admits the same gain.
On $s\le m/256$ all $p_j\ge1/2$. For $x=\alpha s$ and
$\Delta=v-x$, under the success-weighted subset law the exact bias
identity and negative inclusion covariances give
\[
 0\le\E\Delta\le CHx/m,
 \qquad \operatorname{Var}\Delta\le CHx/m.
\]
For the variance, subtract $\min_j z_j$ from all coefficients; the
subset size is fixed, and the remaining nonnegative coefficients
have squared sum at most $H\sum_jz_j$.
The singular Taylor estimate for $P=Q_1Q_2$ is
\[
 |P(v)-P(x)-P'(x)\Delta|
 \le CBd e^{\theta\max(x,v)}\Delta^2/x,
 \quad |P'(x)|\le CB\sqrt{d/x}e^{\theta x}.
\]
At $x=0$ all deficits vanish. Thus, after multiplication by
$\overline A\le e^{-\alpha s}$ and normalization, the error is at most
\[
 CB\frac d m\,m\E_w\left[
  \bigl(H(1+s)+H^2(1+s)/m\bigr)e^{-as}\right].
\]
Split at $L$, use \eqref{extra} and \eqref{laplace} on the lower
part, and reserve half the exponential on the upper part.
The resulting bound is $CB(d/m)(\delta+e^{-cL})$.
The remaining rows $s>m/256$ have error $CB\operatorname{poly}(m)e^{-c_1m}$
by the crude uniform-subset comparison. Decreasing $c$ absorbs this
into the stated bound, since $L\le m/256$ and $d\ge1$.
This proves \eqref{comparison}.

The same Laguerre taper has a positive beta expectation bounded
below absolutely, and factors into Laguerre sums with coefficient
norms $O(d)$ and $O(1)$. Hence its comparison error is
$Cd^2\rho/m$. Choose
\[
 d=\left\lfloor c_2\min\{m,\sqrt{m/\rho}\}\right\rfloor
\]
with a sufficiently small absolute $c_2>0$. The taper then forces
$x_K<2/(\theta d)$. The normalized Christoffel square at this point
has coefficient norm $O(1)$ and beta expectation $O(1/d)$, whence
\[
 \mu(K)\le C/d+Cd\rho/m\le C/d.
\]
As $s_K=O(1/d)$, every $A_B(K)\ge1/2$ for large $m$. Therefore
$w_K\le C/(md)$, exactly the claimed bound.
\end{proof}

For example, a proved profile estimate with
$L=A\log m$ and $\delta\le m^{-\gamma}(\log m)^b$, where
$0<\gamma<1$ and $A$ is sufficiently large, would imply
\[
 K\ge c\,m^{3/2+\gamma/2}(\log m)^{-b/2}.
\]
This last sentence is an implication, not an assertion that such a
quantitative profile has already been established.
\end{document}
